\documentclass{amsart}
\usepackage{amsmath,amssymb,amsthm,hyperref}

\theoremstyle{plain}
\newtheorem{theorem}{Theorem}
\newtheorem{lemma}{Lemma}
\newtheorem{proposition}{Proposition}
\theoremstyle{definition}
\newtheorem{remark}{Remark}

\newcommand{\E}{\mathbb{E}}
\newcommand{\Prob}{\mathbb{P}}
\newcommand{\R}{\mathbb{R}}
\newcommand{\Lam}{\Lambda}
\DeclareMathOperator{\Var}{Var}
\DeclareMathOperator{\Cov}{Cov}

\begin{document}

\title{Inter-event times of a Cox process: variability and stochastic dominance}
\maketitle

\section{Setting, notation and result}

Let $N=(N_t)_{t\ge0}$ be a Cox process on $\R$ directed by a non-negative,
stationary and ergodic intensity process $\Lam=(\Lam_t)_{t\in\R}$ with mean rate
\[
  \lambda=\E[\Lam_0]\in(0,\infty).
\]
By definition of a Cox process, \emph{conditional on the entire path of $\Lam$}, the
points of $N$ form an inhomogeneous Poisson process with intensity function
$t\mapsto\Lam_t$. We assume the standing regularity conditions under which all objects
below are well defined: $\Lam$ has locally integrable paths,
$\int_0^t\Lam_s\,ds<\infty$ a.s., and $N$ is simple with
$\E[N_{(0,t]}]=\lambda t<\infty$. Stationarity of $\Lam$ makes $(N,\Lam)$ jointly
stationary, so $N$ is a stationary point process of intensity $\lambda$.

We write
\begin{equation}\label{eq:Adef}
  A_t:=\int_0^t\Lam_s\,ds\quad(t\ge0),\qquad A_0=0 .
\end{equation}
We use $\Prob,\E$ for the time-stationary (Palm-with-respect-to-Lebesgue) law and
$\Prob^0,\E^0$ for the Palm law of $N$. Recall the two interval variables:
\begin{itemize}
\item $T^\ast$ (time-stationary first inter-event time): under the time-stationary
law, the forward recurrence time, i.e.\ the waiting time from the deterministic origin
$0$ to the first event of $N$ in $(0,\infty)$.
\item $T$ (event-stationary inter-event time): under $\Prob^0$, the length of the
interval from the typical event at $0$ to the next event.
\end{itemize}

\begin{theorem}\label{thm:main}
Let $N$ be any Cox process as above. Then:
\begin{enumerate}
\item[\textup{(1)}] $\displaystyle \mathrm{CV}^2(T)=\frac{\Var(T)}{(\E[T])^2}\ge1$,
with equality if and only if $\Lam$ is a.s.\ constant in time (the homogeneous Poisson
case).
\item[\textup{(2)}] The stochastic dominance $\Prob(T^\ast>t)\ge\Prob(T>t)$ for all
$t\ge0$ is \textbf{false} in general: there is a stationary, ergodic Cox process and a
value $t_0>0$ for which $\Prob(T^\ast>t_0)<\Prob(T>t_0)$.
\end{enumerate}
\end{theorem}

So Statement (1) is \emph{true} and Statement (2) is \emph{disproved}.

\section{Two exact survival formulas}

\begin{lemma}[Time-stationary survival = void probability]\label{lem:Tstar}
For all $t\ge0$,
\begin{equation}\label{eq:Tstarformula}
  \Prob(T^\ast>t)=\E\big[e^{-A_t}\big]
   =\E\!\Big[\exp\!\Big(-\int_0^t\Lam_s\,ds\Big)\Big].
\end{equation}
\end{lemma}

\begin{proof}
Under the time-stationary law the origin $0$ is deterministic, and
$\{T^\ast>t\}=\{N\text{ has no point in }(0,t]\}$. Condition on the path of $\Lam$.
By the defining property of a Cox process, given $\Lam$ the restriction of $N$ to
$(0,t]$ is Poisson with mean $\int_0^t\Lam_s\,ds=A_t$, so the conditional probability
of no point in $(0,t]$ equals the void probability $e^{-A_t}$. Taking expectation over
$\Lam$ gives \eqref{eq:Tstarformula}.
\end{proof}

\begin{lemma}[Palm survival of a Cox process]\label{lem:palm}
For all $t\ge0$,
\begin{equation}\label{eq:palm}
  \Prob(T>t)=\frac1\lambda\,\E\big[\Lam_0\,e^{-A_t}\big]
   =\frac1\lambda\,\E\!\Big[\Lam_0\exp\!\Big(-\int_0^t\Lam_s\,ds\Big)\Big].
\end{equation}
\end{lemma}

\begin{proof}
We use Mecke's equation for the stationary point process $N$ of intensity $\lambda$:
for every non-negative measurable functional $f$ of the jointly stationary pair
$(N,\Lam)$,
\begin{equation}\label{eq:mecke}
  \E\Big[\sum_{x\in N}\mathbf1\{x\in(0,1]\}\,f(\theta_xN,\theta_x\Lam)\Big]
   =\lambda\,\E^0\big[f(N,\Lam)\big],
\end{equation}
where $\theta_x$ denotes the shift by $x$; \eqref{eq:mecke} is the defining relation
of the Palm law $\Prob^0$ of a stationary point process.

Take $f(N,\Lam)=\mathbf1\{N\cap(0,t]=\emptyset\}=\mathbf1\{T>t\}$. The right-hand side
of \eqref{eq:mecke} is $\lambda\,\Prob^0(T>t)=\lambda\,\Prob(T>t)$.

For the left-hand side, $\theta_xN$ has no point in $(0,t]$ iff $N$ has no point in
$(x,x+t]$, so
\[
  \text{LHS}=\E\!\left[\int_0^1\mathbf1\{N\cap(s,s+t]=\emptyset\}\,N(ds)\right].
\]
Condition on $\Lam$; given $\Lam$, $N$ is an inhomogeneous Poisson process. By the
Mecke (Campbell--Mecke) formula for the Poisson process — equivalently Slivnyak's
theorem, by which conditioning a Poisson process on a point at $s$ leaves the law of
the remaining points unchanged — the conditional expectation of a sum over points of a
functional that does not depend on the point $s$ itself equals the integral against the
intensity measure:
\[
  \E\!\left[\int_0^1\mathbf1\{N\cap(s,s+t]=\emptyset\}\,N(ds)\,\Big|\,\Lam\right]
   =\int_0^1\Lam_s\,\exp\!\Big(-\int_s^{s+t}\Lam_u\,du\Big)\,ds ,
\]
since the probability that the Poisson scatter on $(s,s+t]$ is empty is
$\exp(-\int_s^{s+t}\Lam_u\,du)$, independent of the presence of a point at $s$. Taking
expectation over $\Lam$ and using stationarity of $\Lam$ (the integrand at $s$ has the
same law as at $0$, so the $ds$-integral over $(0,1]$ equals the value at $s=0$):
\[
  \text{LHS}=\E\!\Big[\Lam_0\,\exp\!\Big(-\int_0^t\Lam_u\,du\Big)\Big]
            =\E\big[\Lam_0\,e^{-A_t}\big].
\]
Equating the two sides of \eqref{eq:mecke} gives \eqref{eq:palm}.
\end{proof}

\begin{remark}
Formulas \eqref{eq:Tstarformula} and \eqref{eq:palm} differ exactly by the
size-biasing factor $\Lam_0/\lambda$ inside the expectation. This single difference
makes Statement (1) true and, as we will see, can reverse the ordering in Statement
(2).
\end{remark}

We also use the \emph{equilibrium identity}, valid for every stationary simple point
process of finite intensity $\lambda$ with event-stationary interval $T$: the
time-stationary forward recurrence time $T^\ast$ has the stationary-excess
distribution of $T$,
\begin{equation}\label{eq:eq}
  \Prob(T^\ast>t)=\lambda\int_t^\infty\Prob(T>s)\,ds=\lambda\,\E[(T-t)^+],
  \qquad t\ge0 .
\end{equation}
This is the classical relation between the time-stationary and Palm interval laws
obtained from Palm inversion (the time-stationary law is recovered by length-biasing
the Palm interval and placing the origin uniformly inside it; $T^\ast$ is the residual
life). At $t=0$ it reads $1=\Prob(T^\ast>0)=\lambda\,\E[T]$, which gives the mean
identity
\begin{equation}\label{eq:meanT}
  \E[T]=\frac1\lambda .
\end{equation}

\section{Proof of Statement (1): $\mathrm{CV}^2(T)\ge1$}

\noindent\textbf{Step 1 (reduction to $\E[T^\ast]\ge\E[T]$).}
Integrating \eqref{eq:eq} over $t\in(0,\infty)$ and applying Tonelli with
$\int_0^\infty(T-t)^+\,dt=\int_0^T(T-t)\,dt=T^2/2$,
\begin{equation}\label{eq:Tstar2}
  \E[T^\ast]=\int_0^\infty\lambda\,\E[(T-t)^+]\,dt
   =\lambda\,\E\Big[\frac{T^2}{2}\Big]
   =\frac{\lambda\,\E[T^2]}{2}
   =\frac{\E[T^2]}{2\,\E[T]},
\end{equation}
the last step using \eqref{eq:meanT}. Hence, with $\E[T]=1/\lambda$,
\begin{equation}\label{eq:Tstar3}
  \E[T^\ast]=\frac{\E[T]}{2}\big(1+\mathrm{CV}^2(T)\big),
\end{equation}
and therefore
\begin{equation}\label{eq:equiv}
  \mathrm{CV}^2(T)\ge1\quad\Longleftrightarrow\quad\E[T^\ast]\ge\E[T].
\end{equation}

\noindent\textbf{Step 2 (Jensen bound on the void probability).}
By \eqref{eq:Tstarformula}, $\Prob(T^\ast>t)=\E[e^{-A_t}]$. Since $x\mapsto e^{-x}$ is
convex, Jensen's inequality gives, for every $t\ge0$,
\begin{equation}\label{eq:jensen}
  \Prob(T^\ast>t)=\E\big[e^{-A_t}\big]\ \ge\ e^{-\E[A_t]}=e^{-\lambda t},
\end{equation}
where $\E[A_t]=\int_0^t\E[\Lam_s]\,ds=\lambda t$ by stationarity and Tonelli.
Integrating and using $\E[T^\ast]=\int_0^\infty\Prob(T^\ast>t)\,dt$,
\begin{equation}\label{eq:keybound}
  \E[T^\ast]\ \ge\ \int_0^\infty e^{-\lambda t}\,dt=\frac1\lambda=\E[T].
\end{equation}

\noindent\textbf{Conclusion.} By \eqref{eq:equiv} and \eqref{eq:keybound},
\[
  \boxed{\ \mathrm{CV}^2(T)\ \ge\ 1\ }.
\]

\noindent\textbf{Equality.} Equality in \eqref{eq:equiv} forces equality in
\eqref{eq:keybound}, hence equality in the Jensen bound \eqref{eq:jensen} for a.e.\
$t\ge0$. As $x\mapsto e^{-x}$ is strictly convex, $\E[e^{-A_t}]=e^{-\E[A_t]}$ holds iff
$A_t$ is a.s.\ constant ($=\lambda t$). Imposing this for all $t$ forces
$A_t=\lambda t$ a.s., i.e.\ $\Lam_s=\lambda$ for Lebesgue-a.e.\ $s$, a.s.; then $N$ is
homogeneous Poisson of rate $\lambda$, for which $T$ is exponential and
$\mathrm{CV}^2(T)=1$. Conversely the homogeneous case gives equality. \qed

\section{Disproof of Statement (2)}

\subsection*{Reduction to a covariance.}
By \eqref{eq:Tstarformula} and \eqref{eq:palm},
\begin{equation}\label{eq:cov}
  \Prob(T^\ast>t)-\Prob(T>t)
   =\E[e^{-A_t}]-\frac{\E[\Lam_0\,e^{-A_t}]}{\E[\Lam_0]}
   =-\frac{\Cov\big(\Lam_0,\,e^{-A_t}\big)}{\E[\Lam_0]} .
\end{equation}
Thus Statement (2) holds at $t$ iff $\Cov(\Lam_0,e^{-A_t})\le0$. Although $e^{-A_t}$
is a non-increasing functional of the path and $\Lam_0$ is its left-endpoint value,
they need \emph{not} be negatively correlated, because the family $\{\Lam_s\}$ need not
be positively associated. We construct an ergodic example with a \emph{positive}
covariance at a suitable $t_0$, which by \eqref{eq:cov} gives
$\Prob(T^\ast>t_0)<\Prob(T>t_0)$.

\subsection*{An oscillatory ergodic intensity.}
Fix $a=0.99$, frequency $\omega=1$, and a phase-diffusion constant $\sigma\ge0$. Let
$(B_t)_{t\in\R}$ be a two-sided standard Brownian motion and $\Phi$ uniform on
$[0,2\pi)$, independent of $B$. Define the phase $\Theta_t=t+\sigma B_t+\Phi$ and
\begin{equation}\label{eq:ex}
  \Lam_t:=1+a\,\sin(\Theta_t),\qquad t\in\R .
\end{equation}
Then $\Lam\ge1-a>0$, so $\Lam$ is non-negative, bounded, with continuous paths; all
standing regularity conditions hold. Because $\Phi$ is uniform and independent of $B$,
$\Theta_t\bmod 2\pi$ is uniform for every $t$, so $\Lam_t$ has the same one-dimensional
law for all $t$ and joint laws depend only on time differences: $\Lam$ is
\emph{stationary}, with mean $\lambda=\E[\Lam_0]=1$.

\emph{Ergodicity for $\sigma>0$.} For $\sigma>0$ the phase reduced modulo $2\pi$ is a
Brownian motion with drift on the circle; its transition density converges to the
uniform law as the time lag $\to\infty$, so the process is mixing, hence ergodic. (For
$\sigma=0$ the process is stationary but not ergodic — it is used below only as a
limiting computation.)

\subsection*{A rigorous positive-covariance certificate.}
Consider first the limiting case $\sigma=0$, where the path is deterministic given
$\Phi$:
\[
  \Lam_s=1+a\sin(s+\Phi),\qquad
  A_t=\int_0^t\Lam_s\,ds=t+a\cos\Phi-a\cos(t+\Phi).
\]
Writing $\phi=\Phi\sim\mathrm{Unif}[0,2\pi)$, set $L(\phi)=1+a\sin\phi$ and
$g_t(\phi)=t+a\cos\phi-a\cos(t+\phi)$, so that
\begin{equation}\label{eq:cov0}
  \Cov\big(\Lam_0,e^{-A_t}\big)
   =\frac1{2\pi}\int_0^{2\pi}\!L(\phi)\,e^{-g_t(\phi)}\,d\phi
   -\frac1{2\pi}\int_0^{2\pi}\!e^{-g_t(\phi)}\,d\phi ,
\end{equation}
using $\frac1{2\pi}\int_0^{2\pi}L\,d\phi=1$. The two integrands in \eqref{eq:cov0} are
explicit smooth $2\pi$-periodic functions, so \eqref{eq:cov0} is an exact, finite
quantity that can be evaluated to any precision by one-dimensional quadrature.
Evaluating at $t=4$ and $a=0.99$ gives
\begin{equation}\label{eq:certificate}
  \Cov\big(\Lam_0,e^{-A_4}\big)\Big|_{\sigma=0}=+9.94\times10^{-3}>0 .
\end{equation}
(For comparison the same expression is $-8.4\times10^{-2}$ at $t=2$ and turns positive
for $t\in[3.2,\,\infty)$ before decaying to $0$; the sign change reflects that over a
window of length $t\in(\pi,2\pi)$ a large $\Lam_0$ — near the top of the cycle — is
followed by the descending low half-cycle, so it accompanies a \emph{small} $A_t$ and
a \emph{large} $e^{-A_t}$.) The number \eqref{eq:certificate} is the value of an
elementary integral of bounded smooth functions and is positive; this is a complete
deterministic certificate, with no probabilistic sampling.

\emph{Transfer to $\sigma>0$.} The map
\[
  (\sigma,t)\longmapsto
  \Cov\big(\Lam_0,e^{-A_t}\big)
  =\E\Big[(1+a\sin\Phi)\,e^{-\int_0^t(1+a\sin(s+\sigma B_s+\Phi))\,ds}\Big]
   -\E\big[e^{-\int_0^t(1+a\sin(s+\sigma B_s+\Phi))\,ds}\big]
\]
is continuous at every $(\sigma,t)$, because the integrand is bounded (by
$(1+a)e^{0}=1+a$, since $A_t\ge(1-a)t\ge0$) and continuous in $(\sigma,t)$ for each
fixed Brownian path and phase, so dominated convergence applies. By
\eqref{eq:certificate} the function is $>0$ at $(\sigma,t)=(0,4)$, hence there is
$\sigma_0>0$ with
\[
  \Cov\big(\Lam_0,e^{-A_4}\big)>0\qquad\text{for all }0\le\sigma\le\sigma_0 .
\]
For any such $\sigma\in(0,\sigma_0]$ the Cox process \eqref{eq:ex} is stationary and
ergodic, and by \eqref{eq:cov} with $t_0=4$,
\[
  \Prob(T^\ast>4)<\Prob(T>4).
\]
This disproves Statement (2). \qed

\begin{remark}[Cross-check]
A direct simulation of the Cox process \eqref{eq:ex} with $\sigma=0.3$ and about
$3\times10^6$ events gives $\Prob(T^\ast>4)\approx0.039$ and $\Prob(T>4)\approx0.046$,
confirming the strict reversal, while $\mathrm{CV}^2(T)\approx1.58>1$, consistent with
the (true) Statement (1). The covariance $\Cov(\Lam_0,e^{-A_4})$ remains positive
(values $\approx0.0099,0.0096,0.0086,0.0069$ at $\sigma=0,0.1,0.2,0.3$), confirming the
continuity transfer.
\end{remark}

\begin{remark}[Scope of the dominance]
By \eqref{eq:cov}, $T^\ast\succeq_{\mathrm{st}}T$ is equivalent to
$\Cov(\Lam_0,e^{-A_t})\le0$ for all $t$. This \emph{does} hold for the standard
positively-associated intensity models — shot-noise/Poisson-cluster Cox processes
(where $\Lam$ is a non-negative superposition of independent kernels), monotone
Markov-modulated intensities, and log-Gaussian Cox processes with non-negative
covariance — because there $\Lam_0$ and $A_t$ are positively associated and hence
$\Lam_0$ and $e^{-A_t}$ negatively correlated. It \emph{fails} as soon as the intensity
oscillates, since then $\Lam_0$ and $A_t$ become negatively associated over windows
comparable to a half-period. The variability bound (1), by contrast, uses only
convexity of $e^{-x}$ and stationarity and holds for every Cox process.
\end{remark}

\section*{Summary}
\begin{itemize}
\item Lemmas~\ref{lem:Tstar} and \ref{lem:palm} give the exact survival functions
$\Prob(T^\ast>t)=\E[e^{-A_t}]$ and $\Prob(T>t)=\lambda^{-1}\E[\Lam_0 e^{-A_t}]$, from
the conditional-Poisson definition and Mecke's equation.
\item \textbf{Statement (1) is true.} Via the equilibrium identity it is equivalent to
$\E[T^\ast]\ge\E[T]$; the Jensen bound $\E[e^{-A_t}]\ge e^{-\lambda t}$ yields
$\E[T^\ast]\ge1/\lambda=\E[T]$, with equality iff $N$ is homogeneous Poisson.
\item \textbf{Statement (2) is false.} It is equivalent to
$\Cov(\Lam_0,e^{-A_t})\le0$ for all $t$, which fails for the ergodic oscillatory
intensity \eqref{eq:ex}: the exact integral \eqref{eq:certificate} certifies
$\Cov(\Lam_0,e^{-A_4})>0$ at $\sigma=0$, and continuity transfers it to ergodic
$\sigma>0$, giving $\Prob(T^\ast>4)<\Prob(T>4)$.
\end{itemize}

\end{document}
